\subsubsection{Segment Tree Add + lower bound index}
\begin{lstlisting}[style=mystyle, language=C++]
struct SegmentTree {
	struct Node {
		int val = 0, lazy = 0;
	};

	int n;
	vector<Node> st;

	SegmentTree(int size) {
		n = size;
		while(__builtin_popcount(n) != 1) n++;
		st.resize(2*n);
	}

	void push(int node, int l, int r) {
		if(st[node].lazy == 0) return;
		st[node].val += st[node].lazy;
		if(l != r) {
			st[node*2].lazy += st[node].lazy;
			st[node*2+1].lazy += st[node].lazy;
		}
		st[node].lazy = 0;
	}

	void add_range(int node, int l, int r, int ql, int qr, int val) {
		push(node, l, r);
		if(ql <= l && r <= qr) {
			st[node].lazy = val;
			push(node, l, r);
			return;
		}
		if(r < ql || qr < l) return;
		int mid = l + (r-l)/2;
		add_range(node*2, l, mid, ql, qr, val);
		add_range(node*2+1, mid+1, r, ql, qr, val);
		st[node].val = max(st[node*2].val, st[node*2+1].val);
	}

	int query(int node, int l, int r, int val, int bound) {
		push(node, l, r);
		if(r < bound || st[node].val < val) return -1;
		if(l == r) return l;
		int mid = l + (r-l)/2;
		int left = query(node*2, l, mid, val, bound);
		if(left == -1) return query(node*2+1, mid+1, r, val, bound);
		return left;
	}

	void add(int l, int r, int val) {
		add_range(1, 0, n-1, l, r, val);
	}
	int lower_bound_index(int val, int bound) {
		return query(1, 0, n-1, val, bound);
	}
};
\end{lstlisting}
